\documentclass[output=paper]{langscibook}
\ChapterDOI{10.5281/zenodo.22880274}
\title{\textit{O sinivi potee} – The imported canoe called \textit{potee}}
\author{Paul Morekevan and Enoch Horai Magum}
\abstract{}
\begin{document}
\maketitle\label{chap:24}
\begin{figure}[H]
    \includegraphics[trim=0 0 0 0,clip,width=.9\linewidth]{figures/24_1.jpg}
    \caption{\textit{O potee} -- The imported canoe}
    \label{fig:24_1}
\end{figure}
\begin{xltabular}{\textwidth}{lXX}
    (\ref{ex:text18-1})&\textit{O sinivi potee}&The \textit{potee} canoe\\
    (\ref{ex:text18-2})&\textit{O potee o meho kaku sinivi to tei nana Teapu.}&The \textit{potee} is another type of canoe in the Teop area.\\
    (\ref{ex:text18-3})&\textit{Enaa kahi vahutate nom, havee to pakupaku vaavaha me raara e.}&I am going to tell you how we make it.\\
    (\ref{ex:text18-4})&\textit{Beara rake tea paku o potee, eara toro tagihu roho rasuu tea tara rakerake to naono vai to mataa nana tea paku bono potee.}&When we want to make a \textit{potee}, we must first go into the jungle to look for a tree that is good for making the \textit{potee}.\\
    (\ref{ex:text18-5})&\textit{A kiu no potee sa beera potee haana bona kiu vaa teo sinivi suuna,}&The work of carving the \textit{potee} is not as big as the work on the indigenous canoe,\\
    (\ref{ex:text18-6})&\textit{ei mea peha aba na antee nana tea paku bona.}&that is, a single person can do it.\\
    (\ref{ex:text18-7})&\textit{Me bea peha aba rova koa vurahe to naono, eove repaa tagava gunaha bona.}&And even when just a single person finds a (suitable) tree, he fells it.\footnote{Compare Chapter \ref{chap:7}:\ref{ex:text3-5}ff.}\\
    (\ref{ex:text18-8})&\textit{Be tagava bona, eove repaa rom kurus bono govuhu no naono, be tara voen na mataa nana.}&When he has felled it, he cuts the crown of the tree off, if he sees that (the log) is good.\\
    (\ref{ex:text18-9})&\textit{Beve rom kurusu vakavara bona, eove repaa poana pete bona.}&When he has cut it off, he slices the upper side off the log.\\
	(\ref{ex:text18-10})&\textit{A kiu vaa teo potee a varuruva koa tabae o koma beera.}&The work on the potee canoe is easy, because it has a big inside.\\
	(\ref{ex:text18-11})&\textit{Be poana vakavara bona, eove repaa koro bona bona toraaraa ae a rom.}&When having sliced off its upper part, he carves the inside with an axe and a hoe.\\
	(\ref{ex:text18-12})&\textit{A kiu vaa tea koro koma sinivi a beera,}&The work of carving the inside of the indigenous canoe is big,\\
	(\ref{ex:text18-13})&\textit{me be nagunagu, eove repaa kaokahi bona ore paa vaavuriha.}&and when he gets exhausted, he leaves it to have a rest.\\
	(\ref{ex:text18-14})&\textit{A koro no komano sinivi potee na varuruva koana, sa potee haana bono komano sinivi suuna,}&The carving of the inside of the \textit{potee} canoe is quite easy, it is not like the inside of the indigenous canoe,\\
	(\ref{ex:text18-15})&\textit{tabae na antee koa vonaen ei bean parata o buo kubisi, o kubis beera ae o kubis bebeahu.}&because (the indigenous canoe) is only as wide as the span between your thumb and your middle finger.\footnote{lit. `... because only reaches like this: this (is) when you open the two fingers, the big finger and the long finger.'}\\
	(\ref{ex:text18-16})&\textit{Bean boha voen paa anteu, ean repaa nao vo teo hum vaa tea pisipisi.}&When you think that it is enough, then you go on with the part of carving the pointed ends.\\
	(\ref{ex:text18-17})&\textit{Bean taraha pisi vakavara a peha papana, ean repaa nao pete vo tea meha papana.}&When you have finished planing one side of the pointed end, then you go on with the other side.\\
\end{xltabular}
\begin{xltabular}{\textwidth}{lXX}
	(\ref{ex:text18-18})&\textit{Be kavara, ean repaa koro vahabana o komano sinivi.}&After that, you carve again the inside of the canoe.\\
    (\ref{ex:text18-19})&\textit{Koro, koro, bean boha voen, ``antee vai,'' ean repaa nao pete teo hum vaa tea govagovara o sinivi.}&Carve, carve, (and) when you think, ``Enough now,'' you go on with the part of carving the bow and the stern of the canoe.\\
	(\ref{ex:text18-20})&\textit{Bean mei nom tea govara o sinivi, ean toro verete roho ei.}&Before you carve the bow and the stern of the canoe, you must turn it over.\\
	(\ref{ex:text18-21})&\textit{A kiu bona a beera,}&This work is big,\\
	(\ref{ex:text18-22})&\textit{meo aba o sii bero vareko toro paku bari.}&and more than one person must do it.\\
	(\ref{ex:text18-23})&\textit{Beam verete vakavara e, ean repaa taneo a kiu vaa tea govagovara.}&When you have turned it over, you start the work of carving the bow and the stern.\\
	(\ref{ex:text18-24})&\textit{Govara kaakoo teo ihu no sinivi,}&Carve by moving towards the nose of the canoe,\\
	(\ref{ex:text18-25})&\textit{ean repaa vaatotobin e,}&you make it straight,\\
	(\ref{ex:text18-26})&\textit{ean repaa govara pete o hum vaa teo itao.}&then you continue carving the part at the stern.\\
	(\ref{ex:text18-27})&\textit{Bea govara bona kavara, ean paa vaatotobin e.}&When this carving of stern and bow is finished, you make it straight.\\
	(\ref{ex:text18-28})&\textit{Be kavara, ean repaa verete vahae.}&When it is finished, you turn it over again.\\
	(\ref{ex:text18-29})&\textit{O sinivi paa aeae vaevuru,}&The canoe is already light now,\\
	(\ref{ex:text18-30})&\textit{me ean koa a peha na antee nom tea verete vahae.}&and you can turn it over again just by yourself.\\
	(\ref{ex:text18-31})&\textit{Bean verete vakavara, ean repaa koro vahabana e.}&When you have turned it over, you do some more carving on the inside.\\
	(\ref{ex:text18-32})&\textit{Koro, koro, bean boha voen paa anteu, ean paa antee,}&Carve (and) carve, (and) when you think that it is enough, you stop,\\
	(\ref{ex:text18-33})&\textit{ei ihuana vai o hum vaa tea rahi voosu.}&now (you) wait for the part of pulling (the canoe) home.\\
	(\ref{ex:text18-34})&\textit{Teo bon rahi no potee, ean repaa sue ki to sii aba ore ausu vuan tea rahi voosu e.}&For the day of pulling the potee canoe, you tell some people to assist you to pull it home.\\
	(\ref{ex:text18-35})&\textit{A kiu vai ahiki ta von,}&This work does not cost anything,\\
	(\ref{ex:text18-36})&\textit{evehee ean toro paku a sii tabaan to kahi ma an koa vue o manae to paa nao nau tea rahi sinivi.}&but you should make some food that the young men will eat who have gone there to pull the canoe.\\
	(\ref{ex:text18-37})&\textit{Beori rahi voosu maa bari, ean repaa vaan riori, a menaga, nahu iana bara ta maa meha tabaan bata mee.}&When they have pulled it home, you will give them to eat  manioc dumplings, a pot of fish, and also some other food along with it.\\
	(\ref{ex:text18-38})&\textit{Bea maa taba vai kavara, ean paa paku papahiana komana nom amaa sii vagegeana kiu.}&When these things are finished, you will do the final little pieces of work just by yourself.\\
	(\ref{ex:text18-39})&\textit{Beo naono tei me nana bona maa panapanaa, ean repaa gono maa ta maa vua asita tea asita o sinivi tean.}&When the wood has knotholes, you get some putty nuts to seal your canoe with putty nut glue.\\
\end{xltabular}
\begin{figure}[H]
    \captionsetup[subfigure]{oneside,margin={0cm,0cm}}
    \subfloat[]{\includegraphics[trim=0 18 0 18,clip,width=.27\linewidth]{figures/24_2.jpg}\label{fig:24_2}}\hfil
    \subfloat[]{\includegraphics[trim=0 0 0 0,clip,width=.37\linewidth]{figures/24_3.jpg}\label{fig:24_3}}\hfil
    \subfloat[]{\includegraphics[trim=0 0 0 11,clip,width=.36\linewidth]{figures/24_4.jpg}\label{fig:24_4}}
    \caption{(a) \textit{Asita} -- Putty nut tree, (b) and (c) \textit{Vua asita} -- Putty nut}
\end{figure}
\begin{xltabular}{\textwidth}{lXX}
	(\ref{ex:text18-40})&\textit{A meha maa naono na hiki ni nana bona panapanaa.}&Some trees do not have any knotholes.\\
	(\ref{ex:text18-41})&\textit{Eikou ean repaa pahin paku bata koa amaa tapara.}&And so you straightaway make some \textit{tapara} poles (that prevent the sides from moving inside).\\
	(\ref{ex:text18-42})&\textit{O naono vai to tei me nana bona panapanaa toro vaatei roho beo asita raka roho nana.}&The wood that has knotholes must be left untouched until the glue has dried up.\\
	(\ref{ex:text18-43})&\textit{Bean paku vakavara amaa tapara, ean repaa nao vo teo hum vaa tea pam no sinivi.}&When you have finished fixing the \textit{tapara} poles, you move on to the part of tying up the canoe.\\
\end{xltabular}
\begin{figure}[H]
    \includegraphics[trim=0 0 0 90,clip,width=.8\linewidth]{figures/24_5.jpg}
    \caption{\textit{Tapara} -- Fixing poles inside the canoe}
    \label{fig:24_5}
\end{figure}
\begin{xltabular}{\textwidth}{lXX}    
    (\ref{ex:text18-44})&\textit{Teo potee, o buaku koa o vanai, o peho matana, o meho he itao.}&As for the \textit{potee} canoe, there are only two outrigger booms, one near the bow, and the other near the stern.\\
	(\ref{ex:text18-45})&\textit{Bean pam vakavara nio buo vanai, ean repaa tapa ae pam vagegeana koa nio aga ore paa kavara.}&When you have tied the two outrigger booms (to the hull of the canoe), you finally fix the float with pegs and lashings, and then it is finished.\\
	(\ref{ex:text18-46})&\textit{Ei ean antee vai e tea hae tea vagana ge havee.}&Now you can board it to go fishing or whatever.\\
	(\ref{ex:text18-47})&\textit{O kaku sinivi vai o sinivi vaa Solomon.}&This type of canoe is a canoe from the Solomon Islands.\\
	(\ref{ex:text18-48})&\textit{Evehee a ta Solomon saka vaatei haari bono aga.}&But the Solomon people do not attach a float.\\
	(\ref{ex:text18-49})&\textit{Eve he nomaa voen meara paa vaatei o aga,}&When it came to us like this, we attached the float,\\
	(\ref{ex:text18-50})&\textit{ei kou a tahii teara na pohapoha sana nana.}&because our sea is very rough.\\
	(\ref{ex:text18-51})&\textit{O potee na antee nana tea tatana bona too hata.}&The \textit{potee} can resist bad weather.\\
	(\ref{ex:text18-52})&\textit{Sa manin goo haana,}&It never sinks,\\
	(\ref{ex:text18-53})&\textit{o sinivi suuna he na antee nana tea goo.}&whereas the indigenous canoe can sink.\\
\end{xltabular}
\begin{figure}[H]
    \captionsetup{justification=centering}
    \includegraphics[trim=0 0 0 60,clip,width=.9\linewidth]{figures/24_6.jpg}
    \caption{\textit{Sima Rigamu na hua nana bono potee} -- \\Sima Rigamu is paddling a \textit{potee}}
    \label{fig:24_6}
\end{figure}
\end{document}